\subsection{Gauge Fields, Fermions and Mass Gaps in 6D Brane Worlds --- arXiv:hep-th/0608074}

\textbf{Authors:} S. L. Parameswaran, S. Randjbar-Daemi, A. Salvio

\paragraph{主な主張}
6次元 minimal gauged supergravity の axisymmetric warped brane 背景について gauge field と fermion の KK 方程式を Schrödinger 型に帰着し、物理的 boundary condition の下で exact KK wavefunctions と離散質量スペクトルを求め、codimension-two defect により compactification volume を大きくしても有限の KK mass gap を保ち得ることを示した \cite{Parameswaran:2006SixDBrane}。

\paragraph{新規性}
conical defect と warping が zero-mode localization と massive KK tower の双方に与える効果を解析的に制御し、KK mass gap と内部空間体積が必ずしも単純な逆比例関係にない具体例を与えた点にある。

\paragraph{位置付け}
丸い $S^2$ with flux より一般的な warped/rugby-ball compactification の文献であり、compactification geometry や brane defect が higher KK spectrum を変えること、したがって第一 KK mass だけから内部体積を一意に読むことができない場合があることを示す重要な比較対象である。